% ============================================
% Johns Hopkins Letter Template
% Assistant Professor Cover Letter – City College of New York (CUNY)
% ============================================

\documentclass[11pt,letterpaper]{letter}

\usepackage{graphicx}
\usepackage{eso-pic}
\usepackage{geometry}
\usepackage{microtype}
\usepackage[hidelinks]{hyperref}

% ----- Page Setup -----
\geometry{left=25mm,right=25mm,top=25mm,bottom=25mm}

% ----- Logo Placement (first page only) -----
\newcommand{\PutJHULogoFG}[1]{%
  \AddToShipoutPictureFG*{%
    \AtPageUpperLeft{%
      \hspace{10mm}%
      \raisebox{-38mm}{%
        \includegraphics[width=6cm]{#1}%
      }%
    }%
  }%
}

% ----- Sender Information -----
\signature{Decory Edwards}
\address{Department of Economics \\ Johns Hopkins University \\ Baltimore, MD 21218 \\ \texttt{dedwar65@jh.edu}}

\begin{document}
\PutJHULogoFG{Images/jhulogo.png}

\begin{letter}{Search Committee \\
Department of Economics and Business \\
Colin Powell School for Civic and Global Leadership \\
City College of New York (CUNY) \\
160 Convent Avenue \\
New York, NY 10031 \\ USA}

\opening{Dear Members of the Search Committee,}

I am writing to express my interest in the tenure-track Assistant Professor position in Finance/Data Science in the Department of Economics and Business at the Colin Powell School for Civic and Global Leadership at The City College of New York. I am a Ph.D. candidate in Economics at Johns Hopkins University, advised by Professors Christopher D.~Carroll, Nicholas W.~Papageorge, and Stelios Fourakis, and I expect to receive my degree in May 2026. My research fields are macroeconomics, wealth and income dynamics, and computational economics.

My research agenda focuses on understanding the determinants of wealth inequality and the mechanisms through which heterogeneity in returns to wealth shapes aggregate outcomes. In my job-market paper, I estimate the degree of heterogeneity in individual portfolio returns using micro-data from the Survey of Consumer Finances and simulate a structural model in which households face idiosyncratic return risk. I show that allowing for persistent heterogeneity in returns substantially improves the model’s ability to match the empirical distribution of wealth, offering a tractable way to connect micro-level return dispersion to macroeconomic inequality.

In a second project, I use the Health and Retirement Study to examine the relationship between interpersonal trust and portfolio performance. Building on the literature that finds a hump-shaped relationship between trust and income, I show that a similar relationship holds for individual returns to wealth measured in the HRS data, highlighting a behavioral channel through which social attitudes shape saving and investment outcomes. This work also provides new estimates of the persistent component of returns to net worth, which motivate the calibration of heterogeneity in my structural model.

CCNY’s commitment to high-quality research, quantitative rigor, and public service—alongside its mission to serve a uniquely diverse, ambitious, and socially engaged student body—makes the Department of Economics and Business an exciting environment for my work. I am particularly drawn to the position’s focus on finance, data science, financial economics, and quantitative analytics, all of which complement my computational and empirical background in macroeconomics and portfolio behavior. The opportunity to contribute to undergraduate and graduate instruction, supervise MA students, and collaborate across the Colin Powell School’s interdisciplinary programs fits well with my interest in applied data science, financial modeling, and evidence-based economic research.

\textbf{Teaching.} My teaching experience centers on undergraduate macroeconomics and an upper-level macro policy seminar, where I have led weekly discussion sections and held regular office hours for classes ranging from 16 to 40 students. My approach is student-centered: I aim to help students “convince themselves” that they have moved from confusion to understanding by holding structured office-hour coaching that builds trust and confidence. At CCNY, I am prepared to teach \textit{Finance}, \textit{Financial Economics}, \textit{Data Science for Economics}, \textit{Financial Analytics}, \textit{Quantitative Methods}, and related courses at both the undergraduate and master’s levels, and I welcome the opportunity to supervise MA research projects.

I believe I would be a strong fit because my computational and empirical toolkit allows me to (i) teach quantitative and finance-related subjects with rigor and clarity, (ii) integrate reproducible workflows and real data sources into analytics and financial modeling courses, and (iii) mentor students effectively in a diverse, mission-driven institution committed to public value and upward mobility. I am enthusiastic about the opportunity to contribute to CCNY’s teaching, research, and civic mission through the Colin Powell School.

I would be honored to join the Department of Economics and Business at The City College of New York. Thank you for your consideration; I welcome the opportunity to discuss my fit with your department at your convenience.

\closing{Sincerely,}

\end{letter}
\end{document}
